\section{Introducción}

La conversión eficiente de niveles de tensión en los sistemas de distribución eléctrica descansa principalmente sobre el funcionamiento del transformador monofásico. Este dispositivo electromagnético basa su operación en el principio de inducción mutua fundamentado por las leyes de Faraday y Lenz para transferir energía entre sus devanados principal y secundario a través de un núcleo magnético común. Aunque teóricamente se asume un comportamiento ideal en esta transformación, en la práctica el voltaje disponible en el secundario está condicionado tanto por las pérdidas internas de la máquina como por la naturaleza del consumo acoplado a la salida.

Al operar con cargas lineales, como las de carácter resistivo puro, la corriente demandada conserva una forma de onda sinusoidal y proporcional a la tensión, lo que permite un análisis predecible mediante el circuito equivalente del dispositivo. Por el contrario, la presencia de cargas no lineales tales como los rectificadores de onda completa con filtrado altera considerablemente esta dinámica al demandar corriente en forma de pulsos. Este fenómeno introduce distorsión armónica, altera la forma de onda de la corriente y degrada la relación entre la potencia activa y la potencia aparente. Asimismo, dicho comportamiento incrementa las pérdidas por efecto Joule y en el núcleo, elevando la temperatura de operación y afectando el rendimiento global de la máquina.

En este contexto, la presente práctica tiene como propósito analizar experimentalmente el comportamiento de un transformador monofásico sometido a ambos escenarios de carga. Se determinará y contrastará la curva de regulación de tensión para una carga lineal frente al modelo teórico, para posteriormente evaluar el impacto de las cargas no lineales en la distorsión de la corriente y la eficiencia de la transferencia de potencia.